% Options for packages loaded elsewhere
% Options for packages loaded elsewhere
\PassOptionsToPackage{unicode}{hyperref}
\PassOptionsToPackage{hyphens}{url}
%
\documentclass[
  11pt,
  letterpaper,
  DIV=11,
  numbers=noendperiod,
  twoside]{scrartcl}
\usepackage{xcolor}
\usepackage[left=1.1in, right=1in, top=0.8in, bottom=0.8in,
paperheight=9.5in, paperwidth=7in, includemp=TRUE, marginparwidth=0in,
marginparsep=0in]{geometry}
\usepackage{amsmath,amssymb}
\setcounter{secnumdepth}{3}
\usepackage{iftex}
\ifPDFTeX
  \usepackage[T1]{fontenc}
  \usepackage[utf8]{inputenc}
  \usepackage{textcomp} % provide euro and other symbols
\else % if luatex or xetex
  \usepackage{unicode-math} % this also loads fontspec
  \defaultfontfeatures{Scale=MatchLowercase}
  \defaultfontfeatures[\rmfamily]{Ligatures=TeX,Scale=1}
\fi
\usepackage{lmodern}
\ifPDFTeX\else
  % xetex/luatex font selection
  \setmathfont[]{Garamond-Math}
\fi
% Use upquote if available, for straight quotes in verbatim environments
\IfFileExists{upquote.sty}{\usepackage{upquote}}{}
\IfFileExists{microtype.sty}{% use microtype if available
  \usepackage[]{microtype}
  \UseMicrotypeSet[protrusion]{basicmath} % disable protrusion for tt fonts
}{}
\usepackage{setspace}
% Make \paragraph and \subparagraph free-standing
\makeatletter
\ifx\paragraph\undefined\else
  \let\oldparagraph\paragraph
  \renewcommand{\paragraph}{
    \@ifstar
      \xxxParagraphStar
      \xxxParagraphNoStar
  }
  \newcommand{\xxxParagraphStar}[1]{\oldparagraph*{#1}\mbox{}}
  \newcommand{\xxxParagraphNoStar}[1]{\oldparagraph{#1}\mbox{}}
\fi
\ifx\subparagraph\undefined\else
  \let\oldsubparagraph\subparagraph
  \renewcommand{\subparagraph}{
    \@ifstar
      \xxxSubParagraphStar
      \xxxSubParagraphNoStar
  }
  \newcommand{\xxxSubParagraphStar}[1]{\oldsubparagraph*{#1}\mbox{}}
  \newcommand{\xxxSubParagraphNoStar}[1]{\oldsubparagraph{#1}\mbox{}}
\fi
\makeatother


\usepackage{longtable,booktabs,array}
\newcounter{none} % for unnumbered tables
\usepackage{calc} % for calculating minipage widths
% Correct order of tables after \paragraph or \subparagraph
\usepackage{etoolbox}
\makeatletter
\patchcmd\longtable{\par}{\if@noskipsec\mbox{}\fi\par}{}{}
\makeatother
% Allow footnotes in longtable head/foot
\IfFileExists{footnotehyper.sty}{\usepackage{footnotehyper}}{\usepackage{footnote}}
\makesavenoteenv{longtable}
\usepackage{graphicx}
\makeatletter
\newsavebox\pandoc@box
\newcommand*\pandocbounded[1]{% scales image to fit in text height/width
  \sbox\pandoc@box{#1}%
  \Gscale@div\@tempa{\textheight}{\dimexpr\ht\pandoc@box+\dp\pandoc@box\relax}%
  \Gscale@div\@tempb{\linewidth}{\wd\pandoc@box}%
  \ifdim\@tempb\p@<\@tempa\p@\let\@tempa\@tempb\fi% select the smaller of both
  \ifdim\@tempa\p@<\p@\scalebox{\@tempa}{\usebox\pandoc@box}%
  \else\usebox{\pandoc@box}%
  \fi%
}
% Set default figure placement to htbp
\def\fps@figure{htbp}
\makeatother


% definitions for citeproc citations
\NewDocumentCommand\citeproctext{}{}
\NewDocumentCommand\citeproc{mm}{%
  \begingroup\def\citeproctext{#2}\cite{#1}\endgroup}
\makeatletter
 % allow citations to break across lines
 \let\@cite@ofmt\@firstofone
 % avoid brackets around text for \cite:
 \def\@biblabel#1{}
 \def\@cite#1#2{{#1\if@tempswa , #2\fi}}
\makeatother
\newlength{\cslhangindent}
\setlength{\cslhangindent}{1.5em}
\newlength{\csllabelwidth}
\setlength{\csllabelwidth}{3em}
\newenvironment{CSLReferences}[2] % #1 hanging-indent, #2 entry-spacing
 {\begin{list}{}{%
  \setlength{\itemindent}{0pt}
  \setlength{\leftmargin}{0pt}
  \setlength{\parsep}{0pt}
  % turn on hanging indent if param 1 is 1
  \ifodd #1
   \setlength{\leftmargin}{\cslhangindent}
   \setlength{\itemindent}{-1\cslhangindent}
  \fi
  % set entry spacing
  \setlength{\itemsep}{#2\baselineskip}}}
 {\end{list}}
\usepackage{calc}
\newcommand{\CSLBlock}[1]{\hfill\break\parbox[t]{\linewidth}{\strut\ignorespaces#1\strut}}
\newcommand{\CSLLeftMargin}[1]{\parbox[t]{\csllabelwidth}{\strut#1\strut}}
\newcommand{\CSLRightInline}[1]{\parbox[t]{\linewidth - \csllabelwidth}{\strut#1\strut}}
\newcommand{\CSLIndent}[1]{\hspace{\cslhangindent}#1}



\setlength{\emergencystretch}{3em} % prevent overfull lines

\providecommand{\tightlist}{%
  \setlength{\itemsep}{0pt}\setlength{\parskip}{0pt}}



 


\setlength\heavyrulewidth{0ex}
\setlength\lightrulewidth{0ex}
\usepackage[automark]{scrlayer-scrpage}
\clearpairofpagestyles
\cehead{
  Claude
  }
\cohead{
  A Counterexample to Arrhenius’s Sixth Impossibility Theorem
  }
\ohead{\bfseries \pagemark}
\cfoot{}
\makeatletter
\newcommand*\NoIndentAfterEnv[1]{%
  \AfterEndEnvironment{#1}{\par\@afterindentfalse\@afterheading}}
\makeatother
\NoIndentAfterEnv{itemize}
\NoIndentAfterEnv{enumerate}
\NoIndentAfterEnv{description}
\NoIndentAfterEnv{quote}
\NoIndentAfterEnv{equation}
\NoIndentAfterEnv{longtable}
\NoIndentAfterEnv{abstract}
\renewenvironment{abstract}
 {\vspace{-1.25cm}
 \quotation\small\noindent\emph{Abstract}:}
 {\endquotation}
\newfontfamily\tfont{EB Garamond}
\addtokomafont{disposition}{\rmfamily}
\addtokomafont{descriptionlabel}{\rmfamily}
\addtokomafont{title}{\normalfont\itshape}
\let\footnoterule\relax

\makeatletter
\renewcommand{\@maketitle}{%
  \newpage
  \null
  \vskip 2em%
  \begin{center}%
  \let \footnote \thanks
    {\itshape\huge\@title \par}%
    \vskip 0.5em%  % Reduced from default
    {\large
      \lineskip 0.3em%  % Reduced from default 0.5em
      \begin{tabular}[t]{c}%
        \@author
      \end{tabular}\par}%
    \vskip 0.5em%  % Reduced from default
    {\large \@date}%
  \end{center}%
  \par
  }
\makeatother
\RequirePackage{lettrine}

\renewenvironment{abstract}
 {\quotation\small\noindent\emph{Abstract}:}
 {\endquotation\vspace{-0.02cm}}

\setmainfont{EB Garamond Math}[
  BoldFont = {EB Garamond SemiBold},
  ItalicFont = {EB Garamond Italic},
  RawFeature = {+smcp},
]

\newfontfamily\scfont{EB Garamond Regular}[RawFeature=+smcp]
\renewcommand{\textsc}[1]{{\scfont #1}}

\renewcommand{\LettrineTextFont}{\scfont}
% Number theorem-like environments straight through (Theorem 1, Theorem 2, ...)
% rather than by section (1.1, 1.2, ...), which is Quarto's LaTeX default.
% Guarded, since Quarto only defines a counter for environments the document uses.
\makeatletter
\AtBeginDocument{%
  \@for\bw@thm:={theorem,lemma,corollary,proposition,conjecture,definition,%
                 example,exercise,refremark,refsolution}\do{%
    \@ifundefined{c@\bw@thm}{}{\expandafter\counterwithout\expandafter{\bw@thm}{section}}%
  }%
}
\makeatother
\KOMAoption{captions}{tableheading}
\makeatletter
\@ifpackageloaded{caption}{}{\usepackage{caption}}
\AtBeginDocument{%
\ifdefined\contentsname
  \renewcommand*\contentsname{Table of contents}
\else
  \newcommand\contentsname{Table of contents}
\fi
\ifdefined\listfigurename
  \renewcommand*\listfigurename{List of Figures}
\else
  \newcommand\listfigurename{List of Figures}
\fi
\ifdefined\listtablename
  \renewcommand*\listtablename{List of Tables}
\else
  \newcommand\listtablename{List of Tables}
\fi
\ifdefined\figurename
  \renewcommand*\figurename{Figure}
\else
  \newcommand\figurename{Figure}
\fi
\ifdefined\tablename
  \renewcommand*\tablename{Table}
\else
  \newcommand\tablename{Table}
\fi
}
\@ifpackageloaded{float}{}{\usepackage{float}}
\floatstyle{ruled}
\@ifundefined{c@chapter}{\newfloat{codelisting}{h}{lop}}{\newfloat{codelisting}{h}{lop}[chapter]}
\floatname{codelisting}{Listing}
\newcommand*\listoflistings{\listof{codelisting}{List of Listings}}
\makeatother
\makeatletter
\makeatother
\makeatletter
\@ifpackageloaded{caption}{}{\usepackage{caption}}
\@ifpackageloaded{subcaption}{}{\usepackage{subcaption}}
\makeatother
\usepackage{bookmark}
\IfFileExists{xurl.sty}{\usepackage{xurl}}{} % add URL line breaks if available
\urlstyle{same}
% fallback for those not using the hyperref driver hyperxmp:
\makeatletter
\@ifundefined{xmpquote}{\newcommand{\xmpquote}[1]{#1}}{}
\makeatother
\hypersetup{
  pdftitle={A Counterexample to Arrhenius's Sixth Impossibility Theorem},
  hidelinks,
  pdfcreator={LaTeX via pandoc}}


\title{A Counterexample to Arrhenius's Sixth Impossibility Theorem}
\author{}
\date{2026}
\begin{document}
\maketitle
\begin{abstract}
Gustaf Arrhenius's sixth impossibility theorem says that no population
axiology satisfies Egalitarian Dominance, General Non-Extreme Priority,
Non-Elitism, Weak Non-Sadism, and the Weak Quality Addition Condition.
Teruji Thomas has noted a gap in the proof, and Elliott Thornley has
restated the last condition with two quantifiers reversed so that the
proof goes through. I show that the gap cannot be closed: the theorem as
stated is false. A complete, transitive axiology on Arrhenius's own
integer welfare levels satisfies all five conditions in their exact
formulations. It ranks populations first by whether their misery exceeds
what their excellent lives can offset, and otherwise by total welfare;
it accepts the Repugnant Conclusion and rejects the Very Repugnant
Conclusion, the combination the theorem was built to rule out. The
repaired theorem is true, and I show what its extra premise buys. The
four remaining conditions force only a background-relative Very
Repugnant Conclusion; for every axiology on which misery can be
compensated by lives barely worth living, Weak Quality Addition in
either form entails the Quality Addition Condition the sixth theorem was
meant to weaken; and the strengthened condition excludes, by
stipulation, axiologies on which the quantity of misery that marginal
lives can never make up for varies with the population it is added to.
Every theorem stated is verified in Lean 4, and the counterexample was
found by a search method described here.
\end{abstract}


\setstretch{1.1}
\section{The theorem}\label{sec-intro}

Population ethics has a body of impossibility theorems that do for the
comparison of differently sized populations what Arrow's theorem did for
voting: each lists conditions on a betterness ordering of populations,
each apparently undeniable, and shows that no ordering satisfies them
all. The problem they formalize is Parfit's
(\citeproc{ref-Parfit1984}{1984}, ch.~17;
\citeproc{ref-Parfit2016}{2016}), that the ways of avoiding the
Repugnant Conclusion each cost something no less repugnant. After Ng
(\citeproc{ref-Ng1989}{1989}) and Blackorby and Donaldson
(\citeproc{ref-BlackorbyDonaldson1991}{1991}), Gustaf Arrhenius proved
six (\citeproc{ref-Arrhenius2000a}{Arrhenius 2000b},
\citeproc{ref-Arrhenius2000b}{2000a},
\citeproc{ref-Arrhenius2003}{2003}, \citeproc{ref-Arrhenius2009}{2009},
\citeproc{ref-Arrhenius2011}{2011}), now the standard reference point
(\citeproc{ref-Greaves2017}{Greaves 2017}), and regards the sixth as
best, because its conditions are logically the weakest. This paper shows
that it is false, and that the repair which makes it true adds exactly
one stipulation: that the quantity of misery which lives barely worth
living can never make up for must be the same whatever population it is
added to. The theorem says:

\begin{quote}
\textbf{The Sixth Impossibility Theorem.} There is no population
axiology which satisfies the Egalitarian Dominance, the General
Non-Extreme Priority, the Non-Elitism, the Weak Non-Sadism, and the Weak
Quality Addition Condition. (\citeproc{ref-Arrhenius2011}{Arrhenius
2011, 9})
\end{quote}

The theorem has been treated as sound. The Oxford Handbook of Population
Ethics takes it as a point of departure
(\citeproc{ref-ArrheniusEtAl2022}{Arrhenius et al. 2022}), and
Arrhenius's own most recent statement, this year, gives the five
conditions in the same form, with Weak Quality Addition beginning ``For
any population X, there is a perfectly equal population with very high
positive welfare, and a very negative welfare level, and a number of
lives at this level'' (\citeproc{ref-Arrhenius2026}{Arrhenius 2026, sec.
2} and n.~10). Its one sustained critic accepts the proof and contests a
background assumption: Erik Carlson argues that the theorem depends on
Finite Fine-Grainedness, that any two welfare levels are joined by
finitely many slight differences, and grants that ``Arrhenius's proof
shows that no population axiology satisfies the exact versions of his
five adequacy conditions, given that the order of welfare levels assumed
in the proof is fine-grained and A-discrete''
(\citeproc{ref-Carlson2022}{Carlson 2022, sec. 1}). Two readers have
noticed that something is wrong with the proof itself. Teruji Thomas, in
an unpublished reconstruction of all six theorems, found that the
published argument establishes only a weaker conclusion than the one
stated, and noted in print that ``the proof of his favoured sixth
theorem contains an error in the derivation of the `Restricted Quality
Addition Condition' (Arrhenius, 2011, Lemma 1.3). As far as I know, one
has to slightly strengthen the premisses of the theorem, in a way
unlikely to cause further controversy''
(\citeproc{ref-Thomas2018}{Thomas 2018, n. 4}; cf.
\citeproc{ref-Thomas2016}{Thomas 2016}). Thornley, following Thomas,
restates Weak Quality Addition with its first two quantifiers reversed,
remarks that ``the Sixth Impossibility Theorem actually requires the
slightly stronger condition stated here,'' and adds that ``the stronger
version remains a compelling adequacy condition''
(\citeproc{ref-Thornley2021}{Thornley 2021, n. 7}).

Thomas's remark implies that the theorem, not only its proof, has to
change; but neither he nor Thornley shows that it does, and the
strengthening has been treated as a formality. It is not. The theorem as
Arrhenius states it is false, and false on the integer-indexed welfare
levels that Carlson concedes to it. There is a complete, transitive,
anonymous welfarist axiology on exactly those levels that satisfies all
five conditions in the exact formulations Arrhenius gives them. I call
it the \emph{capacity view}. It ranks populations first by whether the
misery they contain exceeds what their excellent lives can offset, and
otherwise by total welfare. It accepts the Repugnant Conclusion, as
total utilitarianism does, and rejects the Very Repugnant Conclusion,
which is the combination of attitudes the move from Arrhenius's earlier
theorems to the sixth was designed to make untenable. So the sixth
theorem does not refute the position it was built to refute.

The strengthened condition does make the theorem true; Thomas proved it
in his note, and the proof is machine-checked here. What I add is an
account of what the strengthening buys. In its original form Weak
Quality Addition says that, for any population, there is some quantity
of very bad lives such that adding them together with any number of
lives barely worth living is no better than adding some number of
excellent lives. The strengthened condition fixes the quantity of very
bad lives first and then quantifies over populations. The difference is
whether the amount of misery that marginal lives can never make up for
may depend on how much excellence the population already contains. The
capacity view says it may, and the four other conditions cannot force
otherwise: what they entail is a Very Repugnant Conclusion relative to a
background chosen after the quantity of misery is fixed, and nothing
stronger. Whether the new premise is as ``compelling'' as the old is
therefore the question whether capacity views deserve to be ruled out by
an adequacy condition rather than by argument. The counterexample was
found by a search method (Section~\ref{sec-method}), and every theorem
here is verified in the Lean 4 proof assistant (Section~\ref{sec-lean}).

\section{The framework and the five conditions}\label{sec-conditions}

Arrhenius's framework is this (\citeproc{ref-Arrhenius2011}{Arrhenius
2011, 4--7}). Welfare levels are quasi-ordered, with a neutral level 0;
Arrhenius assumes \emph{Discreteness}, that between any two levels there
are only finitely many others, arguing that nothing turns on it since
one can pass to a discrete subset of a dense order
(\citeproc{ref-Arrhenius2011}{Arrhenius 2011, 6}). Given Discreteness,
levels are indexed by integers, and \(w+1\) is ``slightly higher'' than
\(w\). Following Thomas (\citeproc{ref-Thomas2016}{2016}) I identify a
population with its welfare distribution and write \(n \cdot x\) for
\(n\) lives at level \(x\) and \(P + Q\) for union; Arrhenius's own
exact formulations write \(A \subset W_x\), \(N(A) = n\) for what I
write \(n \cdot x\), and his rider ``other things being equal'' licenses
the identification.\footnote{Every condition is invariant under
  replacing lives by others at the same level, so nothing is lost by
  working with distributions.} An \emph{axiology} is a reflexive,
transitive, not necessarily complete relation ``at least as good as'';
``better than'' is its strict part.

The conditions refer to a ``very high'' level, a ``very low positive''
range \(R(1,y)\) (the levels from 1 to \(y\)), and a ``very negative''
level. In the exact formulations these are quantified over inside each
condition rather than fixed in advance, which is what makes the
formulations exact; I give those (\citeproc{ref-Arrhenius2011}{Arrhenius
2011, 7--9}), since the result depends on them. A range \(R(x,y)\) is,
by Arrhenius's definition, a union of at least three levels
(\citeproc{ref-Arrhenius2011}{Arrhenius 2011, 7} n.~8).

\begin{quote}
\textbf{Egalitarian Dominance (ED).} If \(N(A) = N(B)\), every member of
\(B\) has welfare below \(x\), and \(A = n \cdot x\), then \(A\) is
better than \(B\).
\end{quote}

\begin{quote}
\textbf{General Non-Extreme Priority (GNEP).} For any level \(z\) there
are a positive level \(u\), a positive range \(R(1,y)\) with \(u > y\),
and a number \(n > 0\) such that for any \(x \geq u\), any
\(B \subset R(1,y)\) with \(N(B) = n\), and any population \(E\):
\(\;n \cdot x + 1\cdot z + E\) is at least as good as
\(B + 1 \cdot (z+1) + E\).
\end{quote}

\begin{quote}
\textbf{Non-Elitism (NE).} For any levels \(x, y\) with \(x - 1 > y\)
there is \(n > 0\) such that for any \(D \subset R(y, x)\):
\(\;(n+1) \cdot (x-1) + D\) is at least as good as
\(1 \cdot x + n \cdot y + D\).
\end{quote}

\begin{quote}
\textbf{Weak Non-Sadism (WNS).} There are a level \(x < 0\) and a number
\(n\) such that for any level \(y > 0\), any \(B \subset W_y\), and any
population \(C\): \(\;B + C\) is at least as good as \(n \cdot x + C\).
\end{quote}

\begin{quote}
\textbf{Weak Quality Addition (WQA).} For any population \(X\) there are
a level \(x < 0\), positive ranges \(R(u,v)\) and \(R(1,y)\) with
\(u > y\), and numbers \(n, m > 0\) such that for any \(z \geq u\) and
any \(B \subset R(1,y)\): \(\;n \cdot z + X\) is at least as good as
\(B + m \cdot x + X\).
\end{quote}

Informally: ED says a perfectly equal population is better than any
same-sized population every member of which is worse off. GNEP says that
lowering one person's welfare by a step is never worse than failing to
raise a fixed number of others from very low positive to very high
welfare, whatever else is true of the population. NE says that one
person's small loss can be compensated by enough others' gains to the
level she falls to, when the rest of the population lies between those
levels. WNS says some number of very bad lives is no better an addition
than any number of lives worth living. WQA says that for any population
there is some quantity of very bad lives such that adding them, with any
number of lives barely worth living, is no better than adding some
number of very good lives.

Three features of WQA matter. First, the comparison is of additions to a
\emph{background} \(X\), and the witnesses are chosen after \(X\) is
given. Second, \(B\) ranges over low-range populations of every size, so
the \(m\) bad lives cannot be compensated by any number of marginal
lives. Third, WQA weakens Arrhenius's Quality Addition Condition, which
drops the \(m\) bad lives and says that for any \(X\) some number of
very good lives added to \(X\) is at least as good as any number of
marginal lives added to \(X\). Quality Addition is the denial of a
background-relative Repugnant Conclusion, which total utilitarians
accept; Weak Quality Addition is the denial of a background-relative
Very Repugnant Conclusion, that for any population of excellent lives
there is a better one consisting of some very bad lives together with
enough lives barely worth living. Arrhenius weakened the condition
because avoiding the Very Repugnant Conclusion is, as he and his
coauthors later put it, ``intuitively more compelling'' than avoiding
the Repugnant Conclusion
(\citeproc{ref-ArrheniusBudolfsonSpears2021}{Arrhenius et al. 2021,
119}): the sixth theorem is meant to bind even those who, like Tännsjö
(\citeproc{ref-Tannsjo2002}{2002}), Huemer
(\citeproc{ref-Huemer2008}{2008}) and a growing group
(\citeproc{ref-ZuberEtAl2021}{Zuber et al. 2021};
\citeproc{ref-SpearsBudolfson2021}{Spears and Budolfson 2021}), have
made their peace with the Repugnant Conclusion.

\section{The capacity view}\label{sec-capacity}

Fix a ``very high'' level \(A \geq 4\), so that the very low positive
levels 1, 2, 3 lie below it. For a population \(P\) let \(H(P)\) be the
number of lives at or above \(A\) (the \emph{excellent} lives),
\(T^-(P)\) the total welfare of the lives with negative welfare, and
\(T(P)\) the total welfare. Let \(I(P) = H(P) + T^-(P)\), and for a
\emph{baseline capacity} \(c \geq 0\) let \(J(P) = \min(0, I(P) + c)\).
\(I\) measures misery against excellence, each excellent life able to
offset one unit of negative welfare; \(c\) is misery that any
population, the empty one included, can absorb; \(J\) is zero when
excellence and baseline cover the misery and otherwise equals the
uncovered misery. The \emph{capacity view} ranks populations lexically
by \((J, T)\):

\begin{quote}
\(P\) is at least as good as \(Q\) if and only if either
\(J(P) > J(Q)\), or \(J(P) = J(Q)\) and \(T(P) \geq T(Q)\).
\end{quote}

This is complete and transitive, being the lexicographic order on pairs
of integers, and it is welfarist and anonymous. A population whose
misery exceeds its capacity to redeem it is worse than any whose misery
does not, and among the former the smaller excess is better; among the
latter, total welfare decides. The exchange rate of one excellent life
per unit of misery is a convenience; any positive weight on \(H\) works.
The simplest member of the family has \(c = 0\), and the reader may take
that as the official countermodel.

The view is lexical, hence discontinuous: a unit of uncovered misery
outweighs any quantity of total welfare. With \(c = 0\), a billion lives
at level 3 together with one life at \(-1\) has \(J = -1\) and is worse
than the empty population; with \(c \geq 1\) the baseline absorbs the
unit and the consequence disappears (Section~\ref{sec-repair}). The view
implies the Repugnant Conclusion, since misery-free populations are
ranked by total welfare. It rejects the Very Repugnant Conclusion: for
\(m|z| > c\), the population \(m \cdot z + N \cdot 3\) has \(I + c < 0\)
and is worse than \(n \cdot A\), which has \(J = 0\), for every \(N\).

It has a close relative in the literature. In the paper that notes the
gap, Thomas sketches a ``toy model'' with ``a sufficientarian flavour'':
integer levels, a sufficiency level \(S\), and populations ``ranked, in
the first instance, by the number of very good lives minus the number of
very bad lives,'' with ties broken by total welfare
(\citeproc{ref-Thomas2018}{Thomas 2018, sec. 4}; cf.
\citeproc{ref-Qizilbash2005}{Qizilbash 2005};
\citeproc{ref-Knapp2007}{Knapp 2007}; \citeproc{ref-Parfit2016}{Parfit
2016}). He offers it for another purpose: it rejects the Repugnant
Conclusion and violates Non-Elitism, but only at the threshold \(S\), so
that if \(S\) is vague the violations are borderline. The capacity view
is Thomas's model with two changes: misery is weighted by welfare rather
than counted, and, decisively, the first layer is capped at zero. The
cap is what turns a determinate violation of Non-Elitism into none,
since the one transfer that lowers \(I\) is licensed only in misery-free
populations, where the capped layer is already zero; and the cap is what
makes the view accept the Repugnant Conclusion.

\textbf{Theorem 1.} For every \(A \geq 4\) and \(c \geq 0\), the
capacity view satisfies ED, GNEP, NE, WNS and WQA in the exact
formulations of Section~\ref{sec-conditions}.\footnote{Arrhenius's 2026
  statement gives the informal conditions with ``better than'' in place
  of ``at least as good as,'' noting that the exact versions are the
  weaker ones (\citeproc{ref-Arrhenius2026}{Arrhenius 2026, n. 10}). The
  capacity view satisfies the strict forms too: ED and WQA are strict
  below, GNEP and WNS have \(\Delta T > 0\), and NE is strict once its
  witness is raised from 1 to 2. The machine-checked statements are the
  non-strict ones.}

Each of the first four conditions compares two populations that differ
by a \emph{transfer}: lives are moved between levels or added. Write
\(\Delta I\) and \(\Delta T\) for the change in \(I\) and \(T\) from the
population the condition says is no better to the one it says is at
least as good. If \(\Delta I \geq 0\) and \(\Delta T \geq 0\) the view
agrees with the condition, since \(J\) is monotone in \(I\) and \(c\)
cancels; backgrounds drop out, because \(I\) and \(T\) are additive over
union.

\emph{Egalitarian Dominance.} \(A = n \cdot x\), every member of \(B\)
is below \(x\), \(N(B) = n\). Then \(T(A) = nx > T(B)\). If \(x \geq A\)
then \(H(A) = n \geq H(B)\); if \(x < A\) then \(H(B) = 0 = H(A)\). And
\(T^-(A) \geq T^-(B)\): if \(x < 0\) every member of \(B\) is negative
and below \(x\), so \(T^-(B) = T(B) < nx = T^-(A)\); if \(x \geq 0\)
then \(T^-(A) = 0 \geq T^-(B)\).

\emph{General Non-Extreme Priority.} Given \(z\), choose \(u = A\),
\(y = 3\), \(n = 2\). Two excellent lives raise \(H\) by 2; the person
moved from \(z+1\) to \(z\) lowers \(H\) by 1 if \(z + 1 = A\) and
lowers \(T^-\) by 1 if \(z + 1 \leq 0\); \(B\) contributes nothing to
\(I\). So \(\Delta I \geq 0\), and \(\Delta T \geq 2A - 7 > 0\).

\emph{Non-Elitism.} Given \(x - 1 > y\), choose \(n = 1\): compare
\(1\cdot x + 1 \cdot y + D\) with \(2 \cdot (x-1) + D\),
\(D \subset R(y,x)\). \(\Delta T = x - 2 - y \geq 0\). If
\(x - 1 \geq A\) the two lives at \(x-1\) are excellent,
\(\Delta H \geq 0\), and \(\Delta T^- \geq 0\) since only the life at
\(y\) can be negative and it is raised. If \(x < A\) then
\(\Delta H = 0\) and \(\Delta T^- \geq 0\); the only case needing care
is \(x = 0\), where two lives at \(-1\) replace one at 0 and one at
\(y \leq -2\), and \(\Delta T^- = |y| - 2 \geq 0\). The remaining case
is \(x = A\), where an excellent life is lost: \(\Delta H = -1\). If
\(y < 0\) the life raised from \(y\) gives \(\Delta T^- = |y| \geq 1\),
so \(\Delta I \geq 0\). If \(y \geq 0\) then \(D \subset R(y, A)\)
contains no negative lives, both populations have \(J = 0\), and \(T\)
decides, and \(T\) does not fall. This is the one place where \(I\) can
fall under a licensed transfer, and it falls only in a population whose
\(J\) is already zero and stays zero. The restriction on \(D\) makes
this work; Arrhenius's General Non-Elitism, used in his fifth theorem,
drops it, and the capacity view violates General Non-Elitism, as it
must, since the fifth theorem is valid.

\emph{Weak Non-Sadism.} Choose \(x = -1\), \(n = 1\): compare
\(1 \cdot (-1) + C\) with \(B + C\), \(B\) at some \(y > 0\).
\(\Delta I = H(B) + 1 \geq 1\) and \(\Delta T = T(B) + 1 \geq 1\).

\emph{Weak Quality Addition.} Given \(X\), choose \(x = -1\), \(u = A\),
\(y = 3\), \(n = 1\), and \(m = \max(0, I(X) + c) + 1\). For any
\(z \geq A\) and \(B \subset R(1,3)\): \(I(1 \cdot z + X) = 1 + I(X)\),
while \(I(B + m \cdot (-1) + X) + c = I(X) + c - m < 0\), since \(B\)
contributes nothing to \(I\). So
\(J(B + m\cdot(-1) + X) < J(1 \cdot z + X)\), and \(1 \cdot z + X\) is
better for every \(B\), whatever its size.

The last argument displays the mechanism. WQA lets the quantity of
misery depend on the background, and the capacity view uses exactly that
freedom: it chooses enough misery to exceed the background's capacity,
after which no quantity of marginal lives helps, because marginal lives
contribute to \(T\) and not to \(I\). The other four conditions never
require comparing a population that has exceeded its capacity favourably
with one that has not, because none of the transfers they license lowers
\(I\) except in populations with no misery.

So the Sixth Impossibility Theorem is false, and Arrhenius's argument
fails at exactly the step Thomas identified. In the proof of Lemma 1.3,
step (12) applies WQA to the background \(A_1 \cup X\) using witnesses
the condition supplies for the background \(X\)
(\citeproc{ref-Arrhenius2011}{Arrhenius 2011, 18--19}), where \(A_1\)
consists of excellent lives whose number depends on those very
witnesses. On the capacity view, adding \(A_1\) raises the background's
capacity, and the misery that exceeded the capacity of \(X\) no longer
exceeds that of \(A_1 \cup X\), so in the enlarged background the
marginal lives win.

\section{The repair and what it costs}\label{sec-repair}

In Thomas's notation (\citeproc{ref-Thomas2016}{2016, 11}), the
conclusion Arrhenius intended is

\begin{quote}
\textbf{Very Repugnant Addition (VRA).} For some population \(P\), any
level \(z < 0\), and any \(m, n \geq 1\), there is \(N\) such that
\(P + m \cdot z + N \cdot 3\) is better than \(P + n \cdot A\),
\end{quote}

and the conclusion the proof delivers is

\begin{quote}
\textbf{Very Repugnant Addition* (VRA*).} For any level \(z < 0\) and
any \(m \geq 1\), there is a population \(P\) such that for any
\(n \geq 1\) there is \(N\) with \(P + m \cdot z + N \cdot 3\) better
than \(P + n \cdot A\).
\end{quote}

In VRA* the background may depend on \(z\) and \(m\). Thomas writes that
VRA* ``is slightly weaker than VRA, because \(P\) is allowed to depend
{[}on{]} \(z\) and \(m\). I think that VRA* is not much less `repugnant'
than VRA, but then again it is not easy to understand the difference
between them'' (\citeproc{ref-Thomas2016}{Thomas 2016, 11}). Thornley's
condition, which in Thomas's fixed-threshold setting entails the denial
of VRA*, is:

\begin{quote}
\textbf{Weak Quality Addition* (WQA*).} There are a level \(x < 0\) and
a number \(m > 0\) such that for any population \(X\) there are a
positive level \(u\), a range \(R(1,y)\) with \(u > y\), and \(n > 0\)
such that for any \(z \geq u\) and any \(B \subset R(1,y)\):
\(\;n \cdot z + X\) is at least as good as \(B + m \cdot x + X\).
\end{quote}

Thomas's corrected theorem is that ED, GNEP, NE and WNS entail VRA*, so
that the five conditions with WQA* are jointly unsatisfiable. I have
verified his proof and proved the unsatisfiability directly for
Arrhenius's exact formulations as well (Section~\ref{sec-lean}). The
capacity view violates WQA*: given \(x\) and \(m\), take \(X\) to be
\(m|x|\) lives at level \(A\), so the misery is exactly covered and both
sides have \(J = 0\); total welfare decides, and enough lives at level 1
in \(B\) make the right-hand side better.

Whether VRA* is much less repugnant than VRA is the question whether
WQA* is as compelling as WQA, and the capacity view says what is at
stake. The two conditions agree on every axiology satisfying:

\begin{quote}
\textbf{Compensation.} For any population \(Y\), any level \(z < 0\) and
any \(m\), there is \(K\) such that \(Y + m \cdot z + K \cdot 3\) is at
least as good as \(Y\).
\end{quote}

Compensation says that any finite misery added to any population can be
made up for by enough lives barely worth living. Total utilitarianism
satisfies it, as do critical-level views with critical level below 3
(\citeproc{ref-BlackorbyBossertDonaldson2005}{Blackorby et al. 2005};
\citeproc{ref-Broome2004}{Broome 2004}), prioritarian views, and any
view Archimedean in the trade between misery and marginal lives. For
such views both WQA and WQA* entail Quality Addition:

\textbf{Proposition 2.} If an axiology satisfies Compensation and WQA,
it satisfies Quality Addition in Arrhenius's exact form: for any \(X\)
there are a positive level \(u\), a range \(R(1,y)\) with \(u > y\), and
\(n > 0\) such that \(n \cdot z + X\) is at least as good as \(B + X\)
for every \(z \geq u\) and every \(B \subset R(1,y)\).

\emph{Proof.} Given \(X\), let \(x, u, y, n, m\) be the witnesses WQA
supplies, and take \(z \geq u\) and \(B \subset R(1,y)\). Compensation
applied to \(X + B\) gives \(K\) with \(X + B + m \cdot x + K \cdot 3\)
at least as good as \(X + B\). Since \(y \geq 3\), \(B + K \cdot 3\)
lies in \(R(1,y)\), so WQA gives \(n \cdot z + X\) at least as good as
\(B + K \cdot 3 + m \cdot x + X\). Transitivity finishes. \(\square\)

The same argument works for WQA*, and the converse holds given Weak
Non-Sadism with a positive number of lives.\footnote{Given \(X\), let QA
  supply \(u, y, n\) and WNS supply \(x'\) and \(n' > 0\). For
  \(B \subset R(1,y)\), WNS with one added life at level 1 and
  background \(B + X\) gives \(B + 1 \cdot 1 + X\) at least as good as
  \(B + n' \cdot x' + X\), and QA, since
  \(B + 1 \cdot 1 \subset R(1,y)\), gives \(n \cdot z + X\) at least as
  good as that; so WQA holds with \(x = x'\), \(m = n'\), and since
  these do not depend on \(X\), so does WQA*.} So for a
Compensation-satisfying axiology the weakening of Quality Addition to
Weak Quality Addition, which distinguishes the sixth theorem from the
fourth, buys nothing. With Thomas's fourth theorem, that ED, GNEP, NE
and WNS entail Repugnant Addition (for any \(X\) and \(n\) there is
\(N\) with \(N \cdot 3 + X\) better than \(n \cdot A + X\), which
entails the denial of Quality Addition), this gives, in Thomas's
fixed-threshold setting:

\textbf{Corollary 3.} ED, GNEP, NE, WNS, WQA and Compensation are
jointly unsatisfiable. Given the four conditions the sixth theorem
shares with the fourth, Weak Quality Addition entails that some finite
quantity of misery cannot be made up for by any number of lives barely
worth living.

The weakening therefore has content only for axiologies that violate
Compensation, and, given the other four conditions, \emph{commits} one
to violating it. Those are the axiologies at which the sixth theorem is
aimed: lexical views on which the badness of a very bad life is of a
different order from the goodness of a marginal one, such as the Lexical
Totalism of Carlson (\citeproc{ref-Carlson2022}{2022}), Thomas
(\citeproc{ref-Thomas2018}{2018}) and Thornley
(\citeproc{ref-Thornley2021}{2021}; cf.
\citeproc{ref-Thornley2022}{Thornley 2022}), and the capacity view.
Among them WQA and WQA* come apart, and the capacity view is the shape
of the gap. WQA* says the quantity of misery that marginal lives cannot
make up for can be fixed once, independently of the population it is
added to; WQA says only that for each population there is such a
quantity. The capacity view takes the second option, because it treats
misery as redeemable by excellent lives and irredeemable by marginal
ones: in a population rich in excellent lives any fixed quantity of
misery is covered, and once covered it is subject to the ordinary
totalist trade against marginal lives; in a population poor in excellent
lives the same quantity is uncovered, and no quantity of marginal lives
touches it.

Is that a structure an adequacy condition should rule out? I do not
think the question has an obvious answer, and I want to resist two quick
ones.

The first is that the capacity view is obviously absurd, so WQA* is as
harmless as Thornley says. The view has features nobody would defend as
stated: an arbitrary exchange rate, a sharp threshold, and an origin in
the search for a model rather than for the truth. With \(c = 0\) it also
says that a billion marginal lives with one slightly bad life is worse
than nothing. That consequence is worth separating from the others,
because it is not a commitment of the structure. The baseline removes
it: with \(c \geq 1\) the empty population absorbs a unit of misery,
every condition still holds, and the Very Repugnant Conclusion still
fails for misery exceeding the baseline. What no member of the family
can remove is uncompensable misery at the empty background once the
misery exceeds \(c\), and that is not peculiar to the capacity view. By
Corollary 3, anyone who accepts the five conditions denies Compensation;
Lexical Totalism in Carlson's and Thornley's versions denies it at every
background, since a very negative life has a negative superior component
that no number of barely-positive lives, which add only inferior goods,
can restore: ``adding a combination of very negative welfare lives and
barely positive welfare lives means adding \((r, s)\) with \(r < 0\)''
(\citeproc{ref-Thornley2021}{Thornley 2021, sec. 4}; cf.
\citeproc{ref-Carlson2022}{Carlson 2022, secs. 4--5}). What the baseline
shows is that the five conditions do not force the uncompensability of a
small quantity of misery. They do not force the uncompensability of mild
misery in any quantity either, though every member of the family as
defined has it, since misery is weighted by welfare from level \(-1\)
down: a variant that counts only lives below a very negative level
\(V\), with weight \(\ell - V\), satisfies the same conditions (checked
numerically, not machine-checked) and makes any number of lives at
\(-1\) compensable at every background. What the five conditions force
is only what Corollary 3 says, that some finite misery is uncompensable
by marginal lives. The sharp threshold is likewise an artefact of taking
the simplest potential: a graded version, in which a life's redemptive
weight rises linearly from 0 at level 3 to 1 at level \(A\), also
satisfies the five conditions (checked numerically for \(A \leq 10\);
the machine-checked proof is for the step version). On that version
every life above the barely-worth-living range helps redeem misery, the
more the higher it lies up to \(A\), and the structural idea is not
absurd. It is close to the thought that the Very Repugnant Conclusion is
worse than the Repugnant one: someone who accepts that enough marginal
lives outweigh excellent lives, but denies that they do when misery has
been thrown in, is treating misery as a dimension that marginal lives do
not reach. The capacity view is the simplest way of making that attitude
consistent with the four other conditions, and that it must let the
uncompensable quantity depend on the background is a discovery about the
attitude, not an artefact of the construction.

The second quick answer is that VRA* is as repugnant as VRA, so the
weakening does not matter. Here the capacity view gives a reason to
disagree with Thomas's first instinct, though I should be clear whose
reason it is. VRA* says that for any quantity of misery there is a
background against which adding it together with enough marginal lives
beats adding excellent lives. On the capacity view the background that
does the work is rich enough in excellent lives to cover the misery, and
what VRA* then asserts is that, in a population whose excellent lives
already redeem its misery, enough additional marginal lives beat
additional excellent lives. That is the Repugnant Conclusion again, for
a population with covered misery in it, and someone who accepts the
Repugnant Conclusion has no new reason to find it repugnant. That is how
the capacity view describes its own verdict. A critic of the Very
Repugnant Conclusion who is not a capacity theorist will describe the
same verdict differently: in some sufficiently good world, adding lives
of torment with enough lives barely worth living is better than adding
excellent lives, and the goodness of the surrounding world does not make
the torment less. I do not think the capacity view answers that critic
without presupposing itself. What it has is a diagnosis. The critic's
complaint is against VRA*, which the four conditions entail; so the
critic must reject one of ED, GNEP, NE and WNS, and Weak Quality
Addition in any form will not help, because the complaint is about a
conclusion the four conditions reach without it. The repaired theorem
tells the critic of the Very Repugnant Conclusion that the fight is
elsewhere. The unrepaired theorem told the critic, wrongly, that there
was no consistent position to fight from.

So the corrected sixth theorem is a theorem about background-relative
repugnance. What ED, GNEP, NE and WNS jointly force is VRA*: for each
quantity of misery, some background against which that misery plus
enough marginal lives beats excellent lives. On the capacity view those
backgrounds are the ones that have redeemed their misery, and the
conclusion is the Repugnant Conclusion for such backgrounds; and the
four conditions do not force VRA, as the capacity view shows. Anyone who
wants more must add WQA*, and the only thing WQA* adds to WQA is the
stipulation that the uncompensable quantity of misery be the same for
every background: by Proposition 2 it has no work to do against
Archimedean views, and its content is the exclusion of views like the
capacity view.

Two remarks on scope. The result does not depend on denying Finite
Fine-Grainedness. The known escape, Carlson's non-Archimedean totalism
(\citeproc{ref-Carlson2022}{2022, sec. 5}), Thomas's total lexic
utilitarianism (\citeproc{ref-Thomas2018}{2018, sec. 1}), Thornley's
Lexical Totalism, gives welfare levels a lexical structure so that no
finite sequence of slight steps leads from very high to very low;
Thornley says ``the theorem is only true given Finite Fine-Grainedness''
(\citeproc{ref-Thornley2021}{Thornley 2021, sec. 3}), and is right about
the strengthened theorem. The capacity view lives on integer levels,
where every level is finitely many steps from every other, and its
lexical structure is in the ranking of populations, not of lives, which
the four fine-grainedness-using conditions do not reach. And the view
also satisfies Dominance Addition, Non-Sadism and Mere Addition; it is
pathological only where it was designed to be.

\section{How the counterexample was found}\label{sec-method}

The capacity view was not found by trying candidates. Arrhenius's
conditions other than the Quality conditions are \emph{transfer
conditions}: each says that a population obtained from another by a
specified change is at least as good, perhaps only when the background
satisfies a restriction. An impossibility proof is a chain of licensed
transfers ending at a population the conclusion's denial says is no
better than the start. So look for a quantity \(\Phi\) on populations
that no licensed transfer lowers, and rank populations first by
\(\min(0, \Phi)\) and then by total welfare. If \(\Phi\) is additive
over union, with a weight \(w_\ell\) per level, backgrounds drop out and
each condition is a linear inequality on the weights. GNEP, lowering one
life from \(z\) to \(z-1\) while \(n\) lives rise from a low level \(y\)
to a high level \(x\), requires \(w_{z-1} + n w_x \geq w_z + n w_y\);
NE, lowering one life from \(x\) to \(x-1\) while \(n\) rise from \(y\)
to \(x-1\), requires \((n+1) w_{x-1} \geq w_x + n w_y\); WNS, comparing
\(D\) lives at a positive level \(y'\) with \(m\) at a negative level
\(x'\), requires \(D w_{y'} \geq m w_{x'}\); and ED requires the weights
to be non-decreasing. One transfer lowers every such \(\Phi\):
Non-Elitism at the very high level, raising non-negative lives, removes
an excellent life and adds nothing negative. It is licensed only in
misery-free backgrounds, so it is neutralized by capping the layer at
zero, provided \(w_\ell \geq 0\) for \(\ell \geq 0\); that is why the
layer is \(\min(0,\Phi)\). And refusing a conclusion of the form ``for
all \(N\), \(m \cdot z + N \cdot 3 + P\) beats \(n \cdot A + P\)''
requires \(\Phi(m \cdot z + N \cdot 3) < 0\) for all \(N\), hence
\(w_3 \leq 0\), hence with the cap \(w_0 = \dots = w_3 = 0\) and
\(w_z < 0\).

For fixed witnesses this is a linear feasibility problem in a dozen
variables, settled instantly by a solver; the script accompanies the
paper. For the sixth theorem's conditions, with every witness 1, it
returns \(w_\ell = \ell\) for \(\ell < 0\), \(0\) for
\(0 \leq \ell < A\), and \(1\) for \(\ell \geq A\): \(\Phi = H + T^-\),
the capacity view. For the fifth theorem, which replaces Non-Elitism by
General Non-Elitism and Weak Non-Sadism by Dominance Addition, the
system is infeasible for every small witness, because the unrestricted
background lets the transfer at the very high level occur in populations
with uncovered misery, forcing \(w_A \leq 0\), which the other
conditions propagate into a contradiction. For the first four theorems
the method is uninformative, since their conclusions involve no misery
and an axiology of this shape can refuse a misery-free conclusion only
through \(w_3 < 0\), which the cap forbids; the sixth is the only
theorem whose conclusion involves misery, and in the fifth General
Non-Elitism closes the gap. The method finds countermodels of one shape,
and infeasibility is evidence rather than proof; but when it succeeds it
delivers a candidate whose verification is routine, and when it fails it
shows which condition excludes. Here it delivered the capacity view in
the first run, and the verification took longer than the search. That
the output was a modification of a model already in the literature is
some evidence that it finds natural objects; the modification, capping
the layer at zero, is the one Thomas's model needed.

\section{Machine verification}\label{sec-lean}

The error in the sixth theorem is a quantifier-dependence error, a
witness chosen for one background used for another; it is the kind a
proof assistant catches and referees do not. Social choice has its
impossibility theorems verified and even discovered by machine
(\citeproc{ref-Nipkow2009}{Nipkow 2009}; \citeproc{ref-TangLin2009}{Tang
and Lin 2009}; \citeproc{ref-GeistEndriss2011}{Geist and Endriss 2011};
\citeproc{ref-BrandtGeist2016}{Brandt and Geist 2016}); population
axiology has one mechanized result that I can find
(\citeproc{ref-ParentBenzmuller2024}{Parent and Benzmüller 2024}), and
none of Arrhenius's theorems had been checked. I have formalized the
results of this paper in Lean 4 (\citeproc{ref-MouraUllrich2021}{Moura
and Ullrich 2021}) using only the core library; the source accompanies
the paper. The theorem \texttt{sixth\_theorem\_countermodel} states that
for every \(A \geq 4\) and \(c \geq 0\) the capacity view is complete
and transitive and satisfies the five conditions in their exact
formulations, and violates WQA*; populations are lists of integer
levels, and ``\(B \subset R(1,y)\)'' is the hypothesis that every member
lies between 1 and \(y\). The theorems \texttt{T1} to \texttt{T5} and
\texttt{T6star} verify Thomas's reconstructions of all six theorems, the
last being that ED, GNEP, NE and WNS entail VRA*, with populations as
finitely supported count functions and conditions in Thomas's
fixed-threshold forms; \texttt{prop2} and \texttt{cor3} verify
Proposition 2 and Corollary 3; and \texttt{T6star\_exact} repeats the
corrected theorem for Arrhenius's exact conditions, where the thresholds
are chosen inside the conditions and the low population is an arbitrary
list in the range. Thomas's statements and proofs are correct in every
case. Two fidelity choices are recorded in the source, besides the side
condition \(y \geq 3\) that Arrhenius's definition of a range imposes:
in Egalitarian Dominance the common size is positive; and in Weak
Non-Sadism, where Arrhenius writes ``a number of lives,'' the
countermodel is shown to satisfy the condition with a positive number,
the stronger claim, while the theorems assume only the weaker form. All
proofs are complete, contain no \texttt{sorry}, and depend only on
Lean's three standard axioms; the development is about twelve hundred
lines.

\section{Conclusion}\label{sec-conclusion}

Arrhenius's sixth impossibility theorem is false. The capacity view
satisfies all five of its conditions in their exact formulations, on
integer-indexed welfare levels, while accepting the Repugnant Conclusion
and rejecting the Very Repugnant Conclusion. The theorem becomes true
when Weak Quality Addition is strengthened so that the uncompensable
quantity of misery is fixed before the background, as Thomas and
Thornley propose, and the strengthened theorem is machine-checked along
with Thomas's reconstructions of the other five. But the strengthening
is not free. For every axiology on which misery can be compensated by
lives barely worth living, Weak Quality Addition in either form entails
the Quality Addition Condition the sixth theorem was meant to weaken;
for the axiologies on which it cannot, the strengthened condition rules
out, by stipulation, the view that how much misery a population can
absorb depends on how much excellence it contains. What the four
remaining conditions establish is that for any quantity of misery there
is some background against which the Repugnant Conclusion holds with
that misery thrown in. On the capacity view those are the backgrounds
that have already redeemed it, and the verdict is the Repugnant
Conclusion over again; a critic who finds the verdict repugnant whatever
the background has a quarrel with the four conditions, not with Weak
Quality Addition in either form.

The error survived years of citation before Thomas found it, and a
decade since, because finding a gap in a proof and finding out whether
the claim is true are different tasks. Population axiology has a stock
of impossibility theorems whose premises are each the product of careful
philosophical judgment and whose proofs are each a page of inequalities
that almost nobody re-derives. The judgment cannot be automated. The
inequalities can, and should be.

\section*{References}\label{references}
\addcontentsline{toc}{section}{References}

\protect\phantomsection\label{refs}
\begin{CSLReferences}{1}{1}
\bibitem[\citeproctext]{ref-Arrhenius2000b}
Arrhenius, Gustaf. 2000a. {``An Impossibility Theorem for Welfarist
Axiologies.''} \emph{Economics and Philosophy} 16 (2): 247--66.
\url{https://doi.org/10.1017/S0266267100000249}.

\bibitem[\citeproctext]{ref-Arrhenius2000a}
Arrhenius, Gustaf. 2000b. \emph{Future Generations: A Challenge for
Moral Theory}. Uppsala University.

\bibitem[\citeproctext]{ref-Arrhenius2003}
Arrhenius, Gustaf. 2003. {``The Very Repugnant Conclusion.''} In
\emph{Logic, Law, Morality: Thirteen Essays in Practical Philosophy in
Honour of Lennart {Å}qvist}, edited by Krister Segerberg and Rysiek
Sliwinski, vol. 51. Uppsala Philosophical Studies. Uppsala University.

\bibitem[\citeproctext]{ref-Arrhenius2009}
Arrhenius, Gustaf. 2009. {``One More Axiological Impossibility
Theorem.''} In \emph{Logic, Ethics, and All That Jazz: Essays in Honour
of Jordan Howard Sobel}, edited by Lars-Göran Johansson, Jan Österberg,
and Rysiek Sliwinski, vol. 57. Uppsala Philosophical Studies. Uppsala
University.

\bibitem[\citeproctext]{ref-Arrhenius2011}
Arrhenius, Gustaf. 2011. {``The Impossibility of a Satisfactory
Population Ethics.''} In \emph{Descriptive and Normative Approaches to
Human Behavior}, edited by Ehtibar N. Dzhafarov and Lacey Perry, vol. 3.
Advanced Series on Mathematical Psychology. World Scientific.
\url{https://doi.org/10.1142/9789814368018_0001}.

\bibitem[\citeproctext]{ref-Arrhenius2026}
Arrhenius, Gustaf. 2026. {``Population Ethics: A Challenge to the
Project of Normative Ethics?''} In \emph{The Oxford Handbook of
Normative Ethics}, edited by David Copp, Connie Rosati, and Tina Rulli.
Oxford University Press.
\url{https://doi.org/10.1093/9780197510926.003.0042}.

\bibitem[\citeproctext]{ref-ArrheniusBudolfsonSpears2021}
Arrhenius, Gustaf, Mark Budolfson, and Dean Spears. 2021. {``Does
Climate Change Policy Depend Importantly on Population Ethics?''} In
\emph{Philosophy and Climate Change}, edited by Mark Budolfson, Tristram
McPherson, and David Plunkett. Oxford University Press.
\url{https://doi.org/10.1093/oso/9780198796282.003.0006}.

\bibitem[\citeproctext]{ref-ArrheniusEtAl2022}
Arrhenius, Gustaf, Krister Bykvist, Tim Campbell, and Elizabeth
Finneron-Burns, eds. 2022. \emph{The Oxford Handbook of Population
Ethics}. Oxford University Press.
\url{https://doi.org/10.1093/oxfordhb/9780190907686.001.0001}.

\bibitem[\citeproctext]{ref-BlackorbyBossertDonaldson2005}
Blackorby, Charles, Walter Bossert, and David Donaldson. 2005.
\emph{Population Issues in Social Choice Theory, Welfare Economics, and
Ethics}. Cambridge University Press.
\url{https://doi.org/10.1017/CCOL0521825512}.

\bibitem[\citeproctext]{ref-BlackorbyDonaldson1991}
Blackorby, Charles, and David Donaldson. 1991. {``Normative Population
Theory: A Comment.''} \emph{Social Choice and Welfare} 8 (3).
\url{https://doi.org/10.1007/BF00177664}.

\bibitem[\citeproctext]{ref-BrandtGeist2016}
Brandt, Felix, and Christian Geist. 2016. {``Finding Strategyproof
Social Choice Functions via {SAT} Solving.''} \emph{Journal of
Artificial Intelligence Research} 55: 565--602.
\url{https://doi.org/10.1613/jair.4959}.

\bibitem[\citeproctext]{ref-Broome2004}
Broome, John. 2004. \emph{Weighing Lives}. Oxford University Press.
\url{https://doi.org/10.1093/019924376X.001.0001}.

\bibitem[\citeproctext]{ref-Carlson2022}
Carlson, Erik. 2022. {``On Some Impossibility Theorems in Population
Ethics.''} In \emph{The Oxford Handbook of Population Ethics}, edited by
Gustaf Arrhenius, Krister Bykvist, Tim Campbell, and Elizabeth
Finneron-Burns. Oxford University Press.
\url{https://doi.org/10.1093/oxfordhb/9780190907686.013.14}.

\bibitem[\citeproctext]{ref-GeistEndriss2011}
Geist, Christian, and Ulle Endriss. 2011. {``Automated Search for
Impossibility Theorems in Social Choice Theory: Ranking Sets of
Objects.''} \emph{Journal of Artificial Intelligence Research} 40:
143--74. \url{https://doi.org/10.1613/jair.3126}.

\bibitem[\citeproctext]{ref-Greaves2017}
Greaves, Hilary. 2017. {``Population Axiology.''} \emph{Philosophy
Compass} 12 (11): e12442. \url{https://doi.org/10.1111/phc3.12442}.

\bibitem[\citeproctext]{ref-Huemer2008}
Huemer, Michael. 2008. {``In Defence of Repugnance.''} \emph{Mind} 117
(468): 899--933. \url{https://doi.org/10.1093/mind/fzn079}.

\bibitem[\citeproctext]{ref-Knapp2007}
Knapp, Christopher. 2007. {``Trading Quality for Quantity.''}
\emph{Journal of Philosophical Research} 32: 211--33.
\url{https://doi.org/10.5840/jpr20073243}.

\bibitem[\citeproctext]{ref-MouraUllrich2021}
Moura, Leonardo de, and Sebastian Ullrich. 2021. {``The {L}ean 4 Theorem
Prover and Programming Language.''} In \emph{Automated Deduction -- CADE
28}, edited by André Platzer and Geoff Sutcliffe, vol. 12699. Lecture
Notes in Computer Science. Springer.
\url{https://doi.org/10.1007/978-3-030-79876-5_37}.

\bibitem[\citeproctext]{ref-Ng1989}
Ng, Yew-Kwang. 1989. {``What Should We Do about Future Generations?
{I}mpossibility of {P}arfit's {T}heory {X}.''} \emph{Economics and
Philosophy} 5 (2): 235--53.
\url{https://doi.org/10.1017/S0266267100002406}.

\bibitem[\citeproctext]{ref-Nipkow2009}
Nipkow, Tobias. 2009. {``Social Choice Theory in {HOL}: {A}rrow and
{G}ibbard-{S}atterthwaite.''} \emph{Journal of Automated Reasoning} 43
(3): 289--304. \url{https://doi.org/10.1007/s10817-009-9147-4}.

\bibitem[\citeproctext]{ref-ParentBenzmuller2024}
Parent, Xavier, and Christoph Benzmüller. 2024. {``Conditional Normative
Reasoning as a Fragment of {HOL}.''} \emph{Journal of Applied
Non-Classical Logics} 34 (4): 561--92.
\url{https://doi.org/10.1080/11663081.2024.2386917}.

\bibitem[\citeproctext]{ref-Parfit1984}
Parfit, Derek. 1984. \emph{Reasons and Persons}. Oxford University
Press.

\bibitem[\citeproctext]{ref-Parfit2016}
Parfit, Derek. 2016. {``Can We Avoid the Repugnant Conclusion?''}
\emph{Theoria} 82 (2): 110--27.
\url{https://doi.org/10.1111/theo.12097}.

\bibitem[\citeproctext]{ref-Qizilbash2005}
Qizilbash, Mozaffar. 2005. {``Transitivity and Vagueness.''}
\emph{Economics and Philosophy} 21 (1): 109--31.
\url{https://doi.org/10.1017/S0266267104000410}.

\bibitem[\citeproctext]{ref-SpearsBudolfson2021}
Spears, Dean, and Mark Budolfson. 2021. {``Repugnant Conclusions.''}
\emph{Social Choice and Welfare} 57 (3): 567--88.
\url{https://doi.org/10.1007/s00355-021-01321-2}.

\bibitem[\citeproctext]{ref-TangLin2009}
Tang, Pingzhong, and Fangzhen Lin. 2009. {``Computer-Aided Proofs of
{A}rrow's and Other Impossibility Theorems.''} \emph{Artificial
Intelligence} 173 (11): 1041--53.
\url{https://doi.org/10.1016/j.artint.2009.02.005}.

\bibitem[\citeproctext]{ref-Tannsjo2002}
Tännsjö, Torbjörn. 2002. {``Why We Ought to Accept the Repugnant
Conclusion.''} \emph{Utilitas} 14 (3): 339--59.
\url{https://doi.org/10.1017/S0953820800003642}.

\bibitem[\citeproctext]{ref-Thomas2016}
Thomas, Teruji. 2016. {``Reconstructing {A}rrhenius's Impossibility
Theorems.''} Unpublished note.
\url{https://users.ox.ac.uk/~mert2060/webfiles/Reconstructing-Arrhenius-for-web.pdf}.

\bibitem[\citeproctext]{ref-Thomas2018}
Thomas, Teruji. 2018. {``Some Possibilities in Population Axiology.''}
\emph{Mind} 127 (507): 807--32.
\url{https://doi.org/10.1093/mind/fzx047}.

\bibitem[\citeproctext]{ref-Thornley2021}
Thornley, Elliott. 2021. {``The Impossibility of a Satisfactory
Population Prospect Axiology (Independently of {F}inite
{F}ine-{G}rainedness).''} \emph{Philosophical Studies} 178 (11):
3671--95. \url{https://doi.org/10.1007/s11098-021-01621-4}.

\bibitem[\citeproctext]{ref-Thornley2022}
Thornley, Elliott. 2022. {``A Dilemma for Lexical and {A}rchimedean
Views in Population Axiology.''} \emph{Economics and Philosophy} 38 (3):
395--415. \url{https://doi.org/10.1017/S0266267121000213}.

\bibitem[\citeproctext]{ref-ZuberEtAl2021}
Zuber, Stéphane, Nikhil Venkatesh, Torbjörn Tännsjö, et al. 2021.
{``What Should We Agree on about the Repugnant Conclusion?''}
\emph{Utilitas} 33 (4): 379--83.
\url{https://doi.org/10.1017/S095382082100011X}.

\end{CSLReferences}




\end{document}
